\documentclass[10pt,a4paper]{amsart}
\usepackage[margin=19mm]{geometry}
\usepackage{amsmath,amssymb,mathtools}
\usepackage[scr=rsfs,cal=euler]{mathalfa}
\usepackage{xfrac}
\usepackage[unicode,hidelinks,bookmarks=false]{hyperref}
\newcommand{\ie}{i.e.\ }
\newcommand{\cf}{cf.\ }
\newcommand{\resp}{respectively\ }
\newcommand{\Subsection}{\subsection}
\providecommand{\nobreakdash}{\nobreak\hbox{-}}
\setlength{\emergencystretch}{2em}
\multlinegap0pt
\usepackage{mathtools}
\usepackage{amssymb}
\usepackage{braket}
\usepackage{bm}
\usepackage{esint}
\usepackage{csquotes}
\usepackage{tikz}
\usepackage{enumerate}
\usepackage{xcolor}
\usepackage{float}
\usepackage{nicefrac}
\usepackage[shortlabels]{enumitem}
\newtheorem{definition}{Definition}
\newtheorem{theorem}{Theorem}
\newtheorem{corollary}{Corollary}
\newtheorem{proposition}{Proposition}
\usepackage{cleveref}
\crefname{definition}{Definition}{Definitions}
\crefname{theorem}{Theorem}{Theorems}
\crefname{corollary}{Corollary}{Corollarys}
\crefname{proposition}{Proposition}{Propositions}
\newcommand{\R}{\mathbb{R}}
\newcommand{\abs}[1]{\left\lvert#1\right\rvert}
\newcommand{\eps}{\varepsilon}
\newcommand{\eqdef}{\ensuremath{\vcentcolon=}}
\title{Positive mass and isoperimetry for continuous metrics with nonnegative scalar curvature}
\author{Gioacchino Antonelli, Mattia Fogagnolo, Stefano Nardulli, Marco Pozzetta}
\date{Source: arXiv:2403.15972v3 (CC BY 4.0)}
\begin{document}
\maketitle

\noindent\emph{Source and licence.} Positive mass and isoperimetry for continuous metrics with nonnegative scalar curvature. Version arXiv:2403.15972v3 (CC BY 4.0), distributed by arXiv under the Creative Commons Attribution 4.0 International licence, http://creativecommons.org/licenses/by/4.0/, as recorded in the arXiv licence metadata for this exact version and shown on the abstract page https://arxiv.org/abs/2403.15972v3. Full licence notice: \texttt{licenses/arxiv-2403.15972-CC-BY-4.0.txt}.\par\smallskip

\noindent\textit{Editorial scope.} Counted unit: the article's theorem thm:CorIsop (source0.tex lines 356--360). The article's complete original proof of it is delivered with it (source0.tex lines 1460--1518, one \texttt{proof} environment of 59 lines, which the article places away from the statement). No mathematics is written, repaired or re-derived: the statement and the proof are the article's own lines, in the article's own notation, and no retained line is substituted or re-flowed.  Retained uncounted as the source-local context the proof needs: lines 272--279; lines 308--310; lines 334--338; lines 489--500; lines 703--707; lines 768--789; lines 1362--1373; lines 1634--1636.  Nothing was omitted from any retained span, so the package carries no omission record.  No retained line contains a citation, so the wrapper carries no bibliography and no bibliography entry.  No retained line contains a figure environment, an image-inclusion command or drawing code, so no image is delivered and the package declares no asset dependency.  Editorial formatting only, in addition: an AMS \texttt{amsart} preamble that restates the article's own declarations and macros used by the retained lines, namely the environments \texttt{definition}, \texttt{theorem}, \texttt{corollary}, \texttt{proposition} on separate counters, together with the standard packages those lines need.  The retained environments print in this document as Definition 1, Theorem 1, Corollary 1, Proposition 1 in order of appearance, whereas the article prints the same items with the numbering of its own full document.  The complete native source of the article as distributed in the arXiv e-print for this exact version is bundled unchanged as \texttt{sources/156539/source0.tex}.
\begin{definition}[$R_g\geq 0$ in the approximate sense]\label{def:NonnegativeScalInApproximate}
    Let $(M,g)$ be a complete $C^0$-Riemannian manifold without boundary, and let $\Omega\subset M$ be an open set. We say that \textit{$R_{g}\geq 0$ in the approximate sense on $\Omega$} if there exist
    smooth complete Riemannian metrics $g_j$ on $M$, such that: 
    \begin{enumerate}
        \item $g_j$ converges to $g$ locally uniformly on $M$;
        \item There exists a sequence $(\varepsilon_j)_{j\in\mathbb N}$ of positive numbers such that $\varepsilon_j\to 0$, and $R_{g_j}\geq -\varepsilon_j$ on $\Omega$.
    \end{enumerate}
\end{definition}

\begin{definition}\label{def:C0locasymptotic}
    Let $K\subset M$ be a compact set (possibly empty) of a $C^0$-Riemannian manifold $(M,g)$. We say that an unbounded connected component $\mathcal{E}$ of $M\setminus K$ is \textit{$C^{0}_{\mathrm{loc}}$-asymptotic to a $C^0$-Riemannian manifold $(N,h)$} if the following holds. For every diverging sequence $\mathcal{E}\ni p_i\to +\infty$ there exists a point $o\in N$ such that $(M,g_i,p_i)\to (N,h,o)$ in the $C^0$-sense, see \cref{def:C0locconvergence}.
\end{definition}

\begin{theorem}
\label{thm:small}
Let $(M,g)$ be a complete 3-dimensional $C^0$-Riemannian manifold without boundary. Let $\Omega \subset M$ be an open subset such that $R_g\ge0$ in the approximate sense on $\Omega$.
Then for any $o \in \Omega$ there exists $r_o>0$ such that for any $r \in (0,r_o]$ there exists $E \subset B_r(o)$ such that $\mathfrak{m}_{\rm QL}(E) \ge 0$.
\end{theorem}

\begin{definition}[$C^0$-convergence]\label{def:C0locconvergence}
We say that pointed $C^0$-Riemannian manifolds \textit{$(M_i,g_i,p_i)$ $C^{0}$-converge to $(M,g,p)$} if the following holds. For every $R,\varepsilon>0$ there exist $i_0:=i_0(R,\varepsilon)$ and an open set $\Omega\subset  M$ such that we have:
\begin{enumerate}
\item $\overline{B}_R(p)\subset  \Omega$;
\item for every $i\geq i_0$ there exists an embedding $F_i:\Omega\to M_i$ such that 
\begin{itemize}
\item $F_i(p)=p_i$;
\item $\overline{B}_R(p_i)\subset  F_i(\Omega)$;
\item $ |(F_i^*g_i-g)_x(v,v)|\leq\varepsilon g_x(v,v)$ for every $x\in \overline{B}_R(p)$ and $v\in T_xM$.
\end{itemize}
\end{enumerate}
\end{definition}

\begin{corollary}\label{cor:Continuity}
Let $(M,g)$ be a $C^0$-Riemannian manifold of dimension $n$ that is $C^0_{\rm loc}$-asymptotic to the $n$-dimensional simply connected complete model $\mathbb H^n_K$ of constant sectional curvature $K\le 0$. Denote by $I$ (resp., $I_K$) the isoperimetric profile of $M$ (resp., $\mathbb{H}^n_K$).

Then $I\le I_K$, and $I$ is locally $\tfrac{n-1}{n}$-H\"{o}lder continuous on $(0,+\infty)$.
\end{corollary}

\begin{theorem}[Asymptotic mass decomposition under $C^0_{\rm loc}$-asymptotic assumptions]\label{thm:MassDecompositionC0}
Let $(M, g)$ be an $n$-dimensional complete $C^0$-Riemannian manifold and assume that $M$ is $C^0_{\rm loc}$-asymptotic to the $n$-dimensional simply connected complete model $\mathbb H^n_K$ of constant sectional curvature $K\le 0$. Fix $o \in M$. Let $V>0$ and let $E_i\subset M$ be a sequence of bounded sets such that $|E_i|=V$ for any $i$ and $\lim_i P(E_i)= I(V)$.

Then, up to subsequence, one of the following alternatives holds true.
\begin{itemize}
    \item The sequence $E_i$ converges in $L^1(M)$ to an isoperimetric set $E$ of volume $V$.

    \item There exist two sequences of radii $R_i, r_i\nearrow+\infty$ with $R_i<r_i$ and a diverging sequence of points $p_i \in M\setminus B_{R_i}(o)$ such that $E_i^c\eqdef E_i \cap B_{R_i}(o)$ converges in $L^1(M)$ to a (possibly empty) isoperimetric set $E$, and $E_i^d\eqdef E_i \cap B_{r_i}(p_i) \setminus B_{R_i}(o)$ converges
    to a ball $B\subset \mathbb H^n_K$ in the sense that
    \[
    \lim_i P(E_i^d) = P(B) , 
        \qquad
    \lim_i |E_i^d| = |B|.
    \]
    Moreover
    \begin{equation*}
        V= |E| + |B|,
        \qquad
        I(V) = P(E) + P(B) .
    \end{equation*}
\end{itemize}
\end{theorem}

\begin{proposition}\label{prop:GrossiQuantoVoglio}
    Let $(M,g)$ be a $3$-dimensional complete $C^0$-Riemannian manifold without boundary, let $K\subset M$ be a compact set, and let $\Omega$ be an unbounded connected component of $M\setminus K$. Assume that $\Omega$ is $C^0_{\mathrm{loc}}$-asymptotic to $\mathbb R^3$ (\cref{def:C0locasymptotic}), and that $R_g\geq 0$ in the approximate sense on $\Omega\setminus K'$ (\cref{def:NonnegativeScalInApproximate}), where $K'\subset M$ is a compact set. Then, there exists a universal constant $\vartheta\in(0,1)$ such that the following holds.
    
    For every $\mathscr{P}>0$, there exists a set of finite perimeter $E\subset\subset \Omega\setminus K'$ such that
    \[
    \vartheta \mathscr{P}\leq P(E)\leq \mathscr{P},
    \]
    and
    \begin{equation}\label{eqn:IlControlloAlContrario}
    |E|\geq \frac{1}{6\sqrt{\pi}}\mathscr{P}^{3/2}\geq \frac{1}{6\sqrt{\pi}}P(E)^{3/2}.
    \end{equation}
\end{proposition}

\begin{theorem}\label{thm:BoundedIsoperimetricResidues}
  Let $(M, g)$ be an $n$-dimensional complete $C^0$-Riemannian manifold that is $C^0_{\mathrm{loc}}$-asymptotic to $\mathbb R^n$. Assume that for some $V > 0$ there exist $E\subset M$ and a nonempty Euclidean ball $B \subset \R^n$ such that $\abs{E} + \abs{B}_{\mathrm{eu}} = V$ and $I(V) = P(E) + P_{\mathrm{eu}}(B)$. Then $E$ has a bounded representative.
\end{theorem}

\begin{theorem}\label{thm:CorIsop}
Let $(M,g)$ be a complete 3-dimensional $C^0$-Riemannian manifold without boundary, and assume that $R_{g}\geq 0$ in the approximate sense on $M\setminus \mathcal{C}$, where $\mathcal{C}$ is a compact set, see \cref{def:NonnegativeScalInApproximate}. Assume in addition that $M$ is $C^0_{\mathrm{loc}}$-asymptotic to $\mathbb R^3$, see \cref{def:C0locasymptotic}.

Then there exists on $M$ a sequence of isoperimetric sets $(E_j)_{j\in\mathbb N}$ such that $|E_j|\to +\infty$ and a sequence of isoperimetric sets $(F_j)_{j\in\mathbb N}$ such that $|F_j|\to 0$. 
\end{theorem}

\begin{proof}[Proof of \cref{thm:CorIsop}]
We deal with the existence of isoperimetric sets of big volume first.
Suppose by contradiction that there exists $v_0 >0$ such that for every $v\in (v_0,+\infty)$ there are no isoperimetric sets of volume $v$ in $M$. We claim that then $I$ is strictly increasing on $(2v_0,+\infty)$.

Let $v\in (2v_0,+\infty)$. By \cref{thm:MassDecompositionC0} there is an isoperimetric set $E\subset M$ with $|E|\leq v_0$ (possibly empty) and a ball $B\subset \mathbb R^3$ with $|B|\geq v-v_0 \geq v/2$ such that 
\[
v=|E|+|B|_{\mathrm{eu}},\qquad I(v)=P(E)+P_{\mathrm{eu}}(B).
\]
Let now $\varepsilon>0$ be such $B_\varepsilon$ is a ball in $\mathbb{R}^3$ concentric to $B$ and with volume $|B_\varepsilon|_{\mathrm{eu}}=|B|_{\mathrm{eu}}-\varepsilon$. Notice that, by approximating $B_\varepsilon\subset \mathbb R^3$ with sets diverging along the manifold, arguing as in the proof of the upper bound for the isoperimetric profile in \cref{cor:Continuity}, we get $I(v-\varepsilon)\leq P(E)+P_{\mathrm{eu}}(B_\varepsilon)$. Thus, taking $\varepsilon\to0$, we find
\[
\frac{I(v)-I(v-\varepsilon)}{\varepsilon} \geq \frac{P_{\mathrm{eu}}(B)-P_{\mathrm{eu}}(B_\varepsilon)}{\varepsilon}
\xrightarrow[\eps\to 0]{}
 2 \left(\frac{4\pi}{3}\right)^{1/3}|B|_{\mathrm{eu}}^{-1/3}
 \geq 2 \left(\frac{4\pi}{3}\right)^{1/3}v^{-1/3}.
\]
    Thus we get that the lower left Dini derivative satisfies
    \[
    \overline{D^-} I(v):=\liminf_{\varepsilon\to 0^+} \frac{I(v)-I(v-\varepsilon)}{\varepsilon} \geq 2 \left(\frac{4\pi}{3}\right)^{1/3}v^{-1/3} 
    >0, \qquad \forall v\in (2v_0,+\infty).
    \]
    Since $I$ is continuous by \cref{cor:Continuity},  the latter implies that $I$ is strictly increasing on $(2v_0,+\infty)$. Thus the sought claim is proved.

    We aim now at showing that there exists an isoperimetric set with volume strictly greater than $2v_0$, thus reaching a contradiction. Fix $v>2v_0$. By \cref{thm:MassDecompositionC0} there is an isoperimetric set $E\subset M$ with $|E|\leq v_0$ and a ball $B\subset \mathbb R^3$ with $|B|\geq v-v_0>0$ such that 
\[
v=|E|+|B|_{\mathrm{eu}},\qquad I(v)=P(E)+P_{\mathrm{eu}}(B).
\]
By \cref{thm:BoundedIsoperimetricResidues} we know that $E$ is bounded. 

We apply \cref{prop:GrossiQuantoVoglio} with $K=\emptyset$ and $K'=\overline{B}_r(o)$ for some ball $B_r(o)$ with $r>1$ such that $E\cup \mathcal{C} \subset B_{r-1}(o)$, where $\mathcal{C}$ is as in the statement. Then there is a set $F\subset M$ such that $F\subset\subset M\setminus K'$, $P(F)\leq P_{\mathrm{eu}}(B)$, and 
\[
|F|\geq \frac{1}{6\sqrt{\pi}}P_{\mathrm{eu}}(B)^{3/2} = |B|_{\mathrm{eu}}.
\]
 Thus the set $E\cup F$ is such that 
\[
|E\cup F| = |E|+|F| \geq |E|+|B|_{\mathrm{eu}} = v.
\]
Hence, since $I$ is strictly increasing on $(2v_0,+\infty)$, one gets
\[
I(|E\cup F|) \geq I(v) = P(E)+P_{\mathrm{eu}}(B),
\]
but at the same time
\[
\begin{split}
I(|E\cup F|) &\leq P(E\cup F)=P(E)+P(F)\leq P(E)+P_{\mathrm{eu}}(B).
\end{split}
\]
Hence all the inequalities in the previous formula are equalities, and then $E\cup F$ is an isoperimetric set with volume $>2v_0$, which is a contradiction. 
\smallskip

We now briefly argue for the existence of isoperimetric sets with small volume. Assume again for a contradiction that there exists $v_0$ such that no isoperimetric set of volume $v$ exists for any $v \in (0, v_0)$. Then, by \cref{thm:MassDecompositionC0}, for every $v\in (0,v_0)$ we have $I(v) = P_{\rm eu}(B_v)$, where $B_v$ is the round ball in $(\R^3, \delta)$ of volume $v$. We can now reach a contradiction as follows. By \cref{thm:small} there is $v'\in (0,v_0)$ and $F$ such that 
\[
|F|=v', \qquad |F|\geq \frac{1}{6\sqrt{\pi}}P(F)^{3/2}.
\]
Thus, since for every $v\in (0,v_0)$ we have $I(v)=P_{\mathrm{eu}}(B_v)$, we can use the information above to conclude that 
\[
P(F)\geq I(v')= \sqrt[3]{36\pi} (v')^{2/3} \geq P(F).
\]
Hence, all the above inequalites are in fact equalities, and $F$ is an isoperimetric set with volume $v'\in (0,v_0)$, reaching a contradiction.
\end{proof}

\end{document}
